%% =====================================================================
%%  T-23-08  The Working Question
%%  -------------------------------------------------------------------
%%  One idea per file. Core is mandatory. Slots in canonical order.
%%  Max 700 words. No layout commands. No filler. New source material
%%  for this entry goes in sources/B7/T-23-08.tex -- never by
%%  rewriting this file (docs/RULES.md R0.1).
%%  =====================================================================
%% @kind: topic
%% @id: T-23-08
%% @title: The Working Question
%% @subtitle: You don't ask, you don't get --- design the question before you need the answer
%% @chapter: c23
%% @order: 80
%% @status: stable
%% @sources: B7
%% @updated: 2026-09-19

\topic{T-23-08}{The Working Question}{You don't ask, you don't get --- design the question before you need the answer}

\begin{core}
B7's craft chapter, drawn on the famous thirty-minute interview that ended in no charges: the answerer starts with the advantage, because they hold the information and you hold only the need for it. Question design is how the questioner levels the field --- one issue per question, no telegraphing what you already know, and the deliberate management of the answerer's \emph{cliff moment}: the line they have decided not to cross in what they will tell you.
\end{core}
\begin{mechanism}
The advantage arithmetic is worth stating plainly: they know what you want to know; you do not know what they want to hide. Every design rule follows from it. One issue per question --- compound questions hand over escape routes, because either half can be answered while the other evaporates. Clean, short, non-accusatory wording --- a question that signals its expected answer teaches the subject what to prepare, so never show your inventory; the famous detectives revealed too much early and spent the half hour on the back foot. Open-then-closed sequencing --- narrative first (which produces testable detail), precise confirmation second. Timing against entrenchment: the subject arrives with a rehearsed story; every fact handed over before it is told helps weld it in place, so the useful questions come early. The cliff moment is the prize --- the answerer has silently fixed how far they will go (\emph{this much, and no more}); the craft is making the step past the line feel survivable: provocative, genuinely neutral questions aimed just past the cliff. And when deception is already known the rule inverts into reserve: never spend what you know to confirm what you suspect --- hold facts back and let the account hang itself.
\end{mechanism}
\begin{conditions}
Full strength in structured one-on-ones where you control the agenda --- interviews, negotiations, the difficult conversation you scheduled. It degrades in group settings (alliances form around answers), against counterparts who recognise interview craft (question met with question), and in intimate relationships, where heavy sequencing reads as ambush and the relational cost can exceed the informational gain.
\end{conditions}
\begin{application}
  \item Before any consequential conversation, write your three questions in order: open, narrower, precise. One issue each. No compounds.
  \item Rank what you know by reveal-order: what you hold back is ammunition; what you show is instruction. Show instruction last.
  \item Ask early. Every question you delay while they talk is a question they will answer later, better rehearsed.
  \item Find the cliff: what have they decided not to say? Aim one question past it, neutrally worded, and let silence do the rest.
  \item After they answer, count to five before speaking (T-23-05). The questioner's silence is the most underrated instrument in the room.
\end{application}
\begin{feedback}
  \item Working: your meetings got shorter and your notes got longer --- the questions are doing the work instead of the arguing.
  \item Working: you notice people rehearsing around your questions now, which is itself information about where the cliffs are.
  \item Failing: they answer the half of the compound question you did not mean. Your design gave them the exit; rebuild the question.
  \item Failing: you revealed your strongest fact in irritation and watched the story reshape itself around it. That was the ammunition budget spent.
\end{feedback}
\begin{failure}
The craft's abuse is the manufactured ambush: questions engineered to trip, reserves deployed to trap, the conversation run as a sting. In any repeatable relationship this converts one win into permanent distrust --- the counterpart who has been interviewed once will never answer casually again, and the information channel narrows to exactly what you can compel. The skill is for questions that must be asked; run constantly, it destroys the ambient honesty it was built to verify.
\end{failure}
\begin{countermeasures}
  \item Notice when someone else's questioning of you has structure: one issue at a time, early, neutral. Craft aimed at you is information about what they suspect.
  \item Answer the question asked, then stop --- the five-second intake on their side is hunting for your second sentence (T-23-05).
  \item In your own defence, ask for questions in writing when stakes are high. Design survives transcription; it withers in improvisation.
  \item Keep your own cliff inventory current: what are you hoping nobody asks? That list is where your leverage leaks.
\end{countermeasures}

\sources{B7}
\seesources
\seealso{T-23-05, T-23-07, T-26-01}
